\documentclass{beamer}
\usepackage{amssymb, latexsym, amsmath, graphics, fullpage, epsfig, amsthm, relsize, pgf, tikz, amsfonts, makeidx, latexsym, ifthen, hyperref, calc}
%\usepackage{eucal} this is a different font than \mathcal{•}
\usetikzlibrary{arrows}
%\newtheorem{theorem}{Theorem}[section]
\newtheorem{remark}{Remark}[section]
%\newtheorem{corollary}{Corollary}
%\newtheorem{lemma}{Lemma}[section]
\newtheorem{assumptions}{Assumptions}[section]
\newtheorem{proposition}{Proposition}
%\newtheorem{definition}{Definition}
\newtheorem{notation}{Notation}
\usepackage{mathrsfs}
%\usepackage[shortlabels]{enumitem}
\usepackage{tikz}
\usepackage{amsmath}
\usetikzlibrary{arrows}
\usetikzlibrary{shadows.blur}
\usepackage{pgfplots}
\usepackage{mwe} % For dummy images
\usepackage{subcaption}
\usepackage{multirow}
\usepackage{dcolumn}
\newcolumntype{2}{D{.}{}{2.0}}
\usepackage{multicol}
\numberwithin{equation}{section}

\newcommand{\lip}{\textup{Lip}_b}
\newcommand{\liph}{\text{Lip}_{\hat\rho}}
\newcommand{\R}{\mathbb{R}}
\newcommand{\E}{\mathbb{E}}
\newcommand{\Norm}[1]{\left\|  #1   \right\|}
\newcommand{\ud}{\ensuremath{\mathrm{d}}}

%\usepackage{lipsum}%% a garbage package you don't need except to create examples.
%\usepackage{fancyhdr}
%\pagestyle{fancy}
%\lhead{\Huge Version A}
%\rhead{\thepage}
%\cfoot{}
%\renewcommand{\headrulewidth}{0pt}



\setbeamersize{text margin left=2mm,text margin right=2mm}

\usepackage{tikz}
\usepackage{pgfplots}
\pgfplotsset{compat=newest}
\usetikzlibrary{decorations.markings}

\allowdisplaybreaks


\defbeamertemplate{footline}{higher page number}
{%
	\hfill%
	\usebeamercolor[fg]{page number in head/foot}%
	\usebeamerfont{page number in head/foot}%
	\insertframenumber\,/\,\inserttotalframenumber\kern1em\vskip20pt%
}
\setbeamertemplate{footline}[higher page number]


%Information to be included in the title page:
\title[About Beamer] %optional
{Interpolating the Stochastic Heat and Wave Equations with Time-independent Noise.}


\author[Eisenberg, Chen] % (optional, for multiple authors)
{Nick~Eisenberg \and Le~Chen}

%\institute[VFU] % (optional)
%{
%	\inst{1}%
%	Faculty of Physics\\
%	UNLV
%	\and
%	\inst{2}%
%	Faculty of Chemistry\\
%	
%}

%\date[Fall 2020] % (optional)
%{Fall 2020}

\begin{document}


\frame{\titlepage}

\begin{frame}[t]
\frametitle{The interpolated stochastic heat and wave equation.}

\begin{equation*}
\begin{cases}
\left(\partial^b_t + \frac{\nu}{2}(-\Delta)^{a/2}\right)\: u(t,x) = I^r_t \left[\sqrt{\theta}\: u(t,x)\: \dot W(x) \right] & \text{$x\in \R^d$, $t>0$}, \\
u(0,\cdot) = 1                                                                                                             & b \in (0,1],               \\
u(0,\cdot) = 1, \quad \partial_t u(0,\cdot) = 0                                                                            & b \in (1,2),
\end{cases}
\end{equation*}


\begin{itemize}
	\item $\partial^b_t$ is the Caputo fractional derivative
	\item $(-\Delta)^{a/2}$ is the fractional Laplacian of order $a \in (0,2]$
	\item $I^r_t$ is the Riemann-Liouville fractional integral of order $r \ge 0$ 
	\item $W = \{ W(\phi) : \phi \in \mathcal{D}(\R^d) \}$ is a centered Gaussian and time-independent noise
\end{itemize}

\end{frame}

\begin{frame}[t]
	\frametitle{Case of $b=a=2$ and $r=0$ and the Riesz kernel noise.}
	
	\begin{equation*}
	\begin{cases}
	\frac{\partial^2 u}{\partial t^2}(t,x) = \Delta u(t,x) + \sqrt{\theta}\: u(t,x)\: \dot W(x) & \text{$x\in \R^d$, $t>0$}, \\
	
	u(0,\cdot) = 1, \quad \partial_t u(0,\cdot) = 0                                                                            & b \in (1,2),
	\end{cases}
	\end{equation*}
	
\begin{theorem}[Balan, Chen, L. and Chen, X.]
	Suppose $\mu(\ud x) = |x|^{-d+\alpha} \ud x$ and that $0 < \alpha < d \le 3$. Let $t_p:=(p-1)^{1/(4-\alpha)}t$. Then the following moment asymptotic holds:
		\[
			\lim_{t_p \to \infty} t_p^{-\frac{4-\alpha}{3-\alpha}} \log \Norm{u(t,x)}_p = \theta^{\frac{1}{3-\alpha}} 2^{\frac{-\alpha}{2(3-\alpha)}} \frac{3- \alpha}{2} \left( \frac{2 \mathcal{M}^{1/2}}{4-\alpha} \right)^{\frac{4-\alpha}{3-\alpha}}.
		\]
In addition by fixing $p\ge 2$, we get that 
	\[
		\lim_{t \to \infty} t^{-\frac{4-\alpha}{3-\alpha}} \log \mathbb{E}|u(t,x)|^p =p(p-1)^{\frac{1}{3-\alpha}} \theta^{\frac{1}{3-\alpha}} 2^{\frac{-\alpha}{2(3-\alpha)}} \frac{3- \alpha}{2} \left( \frac{2 \mathcal{M}^{1/2}}{4-\alpha} \right)^{\frac{4-\alpha}{3-\alpha}}.
	\]
\end{theorem}
\end{frame}

\begin{frame}[t]
	\frametitle{The noise, $W$.}
	We start with a nonnegative and nonnegative definite tempered measure $\Gamma$ with density $\gamma$,
	\[
	\Gamma(\ud x) = \gamma(x) \ud x.
	\]
	Bochner's theorem guarantees the existence of a measure $\mu$, the spectral measure, defined through the following:
	\[
	\int_{\R^d}\Gamma(\ud x) \phi(x) = \int_{\R^d} \mathcal{F}\phi(\xi) \mu(\ud \xi).
	\]
	The covariance for the noise is given as 
	\begin{align*}
	\mathbb{E}(W(\phi)W(\psi)) &= \int_{\R^d} \mathcal{F}\phi(\xi) \overline{\mathcal{F}\psi(\xi)} \mu (\ud \xi)
	\\ &= \int_{\R^{2d}} \phi(x)\psi(y) \gamma(x-y) \ud x \ud y.
	\end{align*}
	When $\mu$ has a density, we denote it as $\varphi(x)$. 
\end{frame}

\begin{frame}[t]
	\frametitle{The noise, $W$.}
\begin{itemize}
	\item Let $k \in \{1,\cdots,d\}$
	\item Partition the $d$-coordinates of $x=(x_1,\cdots,x_n)$ into $k$ distinct groups of size $d_i$ so that $d_1 + \cdots d_k =d$
	\item $x_{(i)} = (x_{i_1},\cdots, x_{i_{d_i}})$ denotes the coordinates of the $i^{th}$ partition.
	\item Choose $\alpha_i \in (0,d_i)$ and denote $\alpha = \sum_{i=1}^k\alpha_i$
\end{itemize}

We consider the following Riesz-type spatial correlation function and its spectral density:
	\[
		\gamma(x) = \prod_{i=1}^k |x_{(i)}|^{-\alpha_i} \quad \varphi(\xi) = \prod_{i=1}^k C_{\alpha_i,d_i} |{\xi}_{(i)}|^{d_i-\alpha_i}.
	\]
	
Note that when $k=1$, the above reduces to the Riesz kernel. 
\end{frame}

\begin{frame}[t]
	\frametitle{The fundamental solution.}
\[
\begin{cases}
\left(\partial^b_t + \frac{\nu}{2}(-\Delta)^{a/2}\right)\: u(t,x) = I^r_t \left[f(t,x)\right] & \text{$x\in \R^d$, $t>0$}, \\
u(0,\cdot) = 1                                                                                                             & b \in (0,1],               \\
u(0,\cdot) = 1, \quad \partial_t u(0,\cdot) = 0                                                                            & b \in (1,2),
\end{cases}
\]

When $f(t,x)$ is a deterministic function, the fundamental solution given through the triple 
	\[
		\{ Z_{a,b,r,\nu,d}(t,x),   Z^*_{a,b,r,\nu,d}(t,x), Y_{a,b,r,\nu,d}(t,x)\}
	\]
where each member is defined explicitly through the Fox-H function. We have for $b \in (0,1]$ and $b \in (1,2)$ respectively that
\[ u(t,x) = 
\begin{cases}
	(Z(t,\cdot)*u(0,\cdot))(x) + (Y\star f)(t,x)
	\\ 	(Z^*(t,\cdot)*u(0,\cdot))(x) + (Z(t,\cdot)*\partial_tu(0,\cdot))(x) + (Y\star f)(t,x).
\end{cases}
\]

	
\end{frame}

\begin{frame}[t]
	\frametitle{Mild solution}
Considering the constant 1 initial condition, it can be shown the solution to the homogeneous equation is equal to $1$. 
\begin{definition}
Let $G(t,x) = Y_{a,b,r,\nu,d}(t,x)$. For $T > 0$, a random field $u=\{u(t,x): t \in (0,T), x \in \R^d \} $ is called a \textit{mild solution} if $G(t-s,x-\cdot)u(s,\cdot)1_{\{s<t\}}$ is Skorohod integrable and the following holds almost surely:
\end{definition}
	\[
		u(t,x) = 1 + \sqrt{\theta} \int_0^t \left(\int_{\R^d} G(t-s,x-y)u(s,y)W(\delta y)\right) \ud s.
	\]
	
\begin{itemize}
	\item \textit{Global Solution:} $\Norm{u(t,x)}_p < \infty$ for any $t>0$ and $x \in \R^d$.
	\item \textit{Local Solution:} If there exists $0 < T_1 \le T_2 < \infty$ such that $\Norm{u(t,x)}_2 < \infty$ when $0<t<T_1$ and $\Norm{u(t,x)}_2 $ D.N.E. for $t> T_2$.
\end{itemize}
\end{frame}

\begin{frame}[t]
	\frametitle{Assumption on $G$}
We need to make an assumption that $G(t,x)$ is nonnegative. This is the case under any of the following:
	\begin{itemize}
		\item $d \ge 1$, $b\in(0,1]$, $a\in(0,2]$, $r\ge0$;
		% \item $d \ge 1$, $b=1$, $a\in (0,2]$, $r=0$ or $r>1$;
		\item $1\le d \le 3$, $1<b<a\le 2$, $r>0$;
		\item $1 \le d \le 3$, $1 < b=a < 2$, $r > \dfrac{d+3}{2}-b$.
	\end{itemize}
\end{frame}

\begin{frame}[t]
	\frametitle{Theorem 1.6 (Chen and Eisenberg)}
Assuming $G(t,x)$ is nonnegative and the noise has spatial correlation $\gamma(x) = \prod_{i=1}^k |x_{(i)}|^{-\alpha_i}$:
	\begin{itemize}
		\item A \textcolor{blue}{\textit{global solution}} exists provided 
			\[
				0 < \alpha < \min\left( \frac{a}{b}[2(b+r) -1], 2a, d \right)
			\]
		\item Otherwise, if 
			\[
					r\in \left[0,1/2\right] \qquad \text{and} \qquad
				0<\alpha = \frac{a}{b}[2(b+r)-1] \le d,
			\]
then a \textcolor{blue}{\textit{local solution}} exists for $t \in (0,T_p))$ where 
	\[
		T_p := \dfrac{\nu^{\alpha/a}}{2 \theta (p-1) \mathcal{M}_a^{(2a-\alpha)/a}}
	\]
and the solution D.N.E. for $t > T_2$.
	\end{itemize}
\end{frame}

\begin{frame}[t]
	\frametitle{Wiener Chaos Expansion:}
If $t=t_{n+1}$ and $x=x_{n+1}$, then a Picard iteration scheme will suggest that for  
	\[
		f_n(x_1, \cdots , x_n; x, t) = \int_0^t\int_0^{t_n}\cdots \int_0^{t_2} \prod_{k=1}^n G(t_{k+1}-t_k,x_{k+1}-x_k) \ud t_1 \cdots \ud t_n,
	\]
the solution can be written as 
	\begin{equation}\label{E:WC}
		u(t,x) = 1 + \sum_{k=1}^\infty \theta^{k/2} I_k(f_k(\cdot,x,t)), \quad (t,x)\in (0,T)\times \R^d.
	\end{equation}
This further suggests 
	\begin{equation}\label{E:2Mom}
		\mathbb{E}(u(t,x)^2) = \sum_{k=0}^\infty \theta^{k} \Norm{f_k(\cdot,x,t)}^2_{\mathcal{H}^{\otimes n}}, \quad (t,x)\in (0,T)\times \R^d.
	\end{equation}
\begin{theorem}
A unique solution exists and \eqref{E:WC} is valid if and only if \eqref{E:2Mom} is finite. 
\end{theorem}
\end{frame}

\begin{frame}[t]
	\frametitle{The moment asymptotic}
We define
	\begin{align*}
			\rho_{\nu,a}(\gamma) &= \sup_{\Norm{f}_{L^2(\R^d)} =1} \int_{\R^d} \left[ \int_{\R^d} \frac{f(x+y)f(y)}{\sqrt{1+\frac{\nu}{2}|x+y|^a } \sqrt{1+\frac{\nu}{2}|y|^a} } \ud y \right]^2 \mu(\ud x)
			\\&= \nu^{-\alpha/a}\mathcal{M}_{a,d}^{2-(\alpha/a)}(\gamma,1)
	\end{align*}
where
	\begin{align*}
		\mathcal{M}_{a,d}(\gamma,\theta) & := \sup_{g \in \mathcal{F}_a} \left\{ \left( \iint_{\R^{2d}} g^2(x)g^2(y) \gamma(x+y)\ud x\ud y \right)^{1/2} - \frac{\theta}{2}\: \mathcal{E}_a(g,g) \right\} \\
		& = \sup_{g \in \mathcal{F}_a} \left\{ \left\langle g^2 * g^2, \gamma \right\rangle_{L^2(\R^d)}^{1/2} - \frac{\theta}{2}\: \mathcal{E}_a(g,g) \right\},
	\end{align*}
and
\begin{gather}
\mathcal{E}_a(g,g) := (2\pi)^{-d} \int_{\R^d} |\xi|^a |\mathcal{F}g(\xi)|^2 \ud \xi \quad \text{and} \\
\mathcal{F}_a      := \left\{f\in L^2(\R^d): \: \Norm{f}_{L^2(\R^d)}=1,\: \mathcal{E}_a(f,f)<\infty \right\}.
\end{gather}
\end{frame}

\begin{frame}[t]
	\frametitle{The moment asymptotic}
\begin{theorem}
	\begin{align*}
		\lim_{t_p \to \infty} &  t_p^{-\beta} \log \Norm{u(t,x)}_p 
		\\& =  \left(\frac{1}{2}\right)\left(\frac{2a}{2a(b+r)- b\alpha} \right)^\beta \\
		& \quad \quad \quad  \times \left(\theta\nu^{-\alpha/a} \mathcal{M}_a^{\frac{2a-\alpha}{a}}\right)^{\frac{a}{2a(b+r)-b\alpha-a}}\left(2(b+r)-\frac{b\alpha}{a}-1\right),
	\end{align*}
where
		\begin{align}
			\label{E:beta-tp}
			\beta := \dfrac{2(b+r)-\dfrac{b\alpha}{a}}{ 2(b+r)-\dfrac{b\alpha}{a} -1 } \qquad \text{and} \qquad
			t_p   := (p-1)^{1-1/\beta} \: t.
		\end{align}
\end{theorem}

\end{frame}

\begin{frame}[t]
	\frametitle{The moment asymptotic}
\begin{theorem} For $p \ge 2$ fixed and $t>0$ fixed respectively, we have that
	\begin{align*}
	1) \lim_{t \to \infty} & t^{-\beta} \log \E\left(|u(t,x)|^p\right) 
	\\&=  p(p-1)^{\frac{1}{2(b+r)-\frac{b\alpha}{a} -1}} \left(\frac{1}{2}\right)\left(\frac{2a}{2a(b+r)- b\alpha} \right)^\beta \\
	& \quad \quad \times \left(\theta\nu^{-\alpha/a} \mathcal{M}_a^{\frac{2a-\alpha}{a}}\right)^{\frac{a}{2a(b+r)-b\alpha-a}}\left(2(b+r)-\frac{b\alpha}{a}-1\right);
	\end{align*}
	\begin{align*}
2) \lim_{p \to \infty} &  p^{-\beta} \log \E\left(|u(t,x)|^p\right) 
\\ &=  t^\beta \left(\frac{1}{2}\right)\left(\frac{2a}{2a(b+r)- b\alpha} \right)^\beta \\
& \quad \quad \times \left(\theta\nu^{-\alpha/a} \mathcal{M}_a^{\frac{2a-\alpha}{a}}\right)^{\frac{a}{2a(b+r)-b\alpha-a}}\left(2(b+r)-\frac{b\alpha}{a}-1\right).
\end{align*}
\end{theorem}
	
	
\end{frame}

\begin{frame}[t]
	\frametitle{Final comments}

\end{frame}
\end{document}